% =====================================================================
%  Modelado para simulacion de un ESC BLDC sensorless con carga
%  agitadora magnetica: ecuaciones, simplificaciones y rango de validez.
%  Compilar con:  latexmk -pdf modelo_simulacion.tex   (o pdflatex x2)
% =====================================================================
\documentclass[11pt,a4paper,oneside]{report}

\usepackage[utf8]{inputenc}
\usepackage[T1]{fontenc}
\usepackage[spanish,es-noshorthands,es-tabla]{babel}
\usepackage{lmodern}
\usepackage{microtype}
\usepackage[margin=2.5cm]{geometry}
\usepackage{amsmath,amssymb,amsthm}
\usepackage{bm}
\usepackage{siunitx}
\sisetup{output-decimal-marker={,},group-separator={.},per-mode=symbol}
\usepackage{booktabs}
\usepackage{array}
\usepackage{enumitem}
\usepackage{xcolor}
\usepackage{tikz}
\usetikzlibrary{arrows.meta,positioning}
\usepackage[hidelinks]{hyperref}

\theoremstyle{definition}
\newtheorem{definicion}{Definición}[chapter]
\newtheorem{ejemplo}{Ejemplo numérico}[chapter]
\theoremstyle{remark}
\newtheorem{nota}{Nota}[chapter]
\newtheorem{hipotesis}{Hipótesis}[chapter]

\newenvironment{idea}%
  {\par\medskip\noindent\begin{minipage}{\linewidth}\hrule\smallskip
   \noindent\textbf{Idea física.}\ }%
  {\smallskip\hrule\end{minipage}\par\medskip}
\newenvironment{lectura}%
  {\par\smallskip\noindent\textbf{Cómo se lee.}\ }%
  {\par\smallskip}

\newcommand{\dd}{\mathrm{d}}
\newcommand{\sgn}{\operatorname{sgn}}
\newcommand{\clamp}{\operatorname{sat}}

\title{\textbf{Modelado de un ESC BLDC sensorless\\ con carga agitadora magnética}\\[0.5em]
\large Modelo mínimo defendible, no linealidades dominantes,\\
linealización y rango de validez}
\author{Proyecto \texttt{esc-bldc-stm32}}
\date{Apuntes de estudio (revisión post-consulta)}

\begin{document}
\hypersetup{pageanchor=false}
\maketitle
\hypersetup{pageanchor=true}
\tableofcontents

% =====================================================================
\chapter{Introducción, alcance y convenciones}
% =====================================================================

\section{Objetivo y criterio de modelado}
Estos apuntes reúnen las ecuaciones necesarias para modelar, linealizar y simular
el sistema completo \emph{motor BLDC (disco duro) + acople magnético no lineal
+ agitador + fluido + sensor de velocidad por cruces por cero + arranque a lazo
abierto + controlador PI discreto}.

Siguiendo el criterio de la cátedra, el cuerpo principal mantiene únicamente un
\textbf{modelo mínimo comprensible y defendible con lápiz y papel}. Se evita
introducir dinámicas o parámetros que no puedan justificarse desde mediciones,
geometría o código real del firmware, para no generar una falsa sensación de
precisión. Las deducciones extendidas y efectos secundarios quedan separados en
el Apéndice~\ref{ap:referencia}.

\paragraph{No linealidades dominantes del problema.} El modelo principal
captura cuatro fenómenos físicos esenciales:
\begin{enumerate}
  \item \textbf{Pérdida de medición a baja velocidad:} por debajo de
        $n_{\text{ciego}}\approx 300\text{--}\SI{400}{rpm}$, la BEMF cae por
        debajo del umbral efectivo del comparador y el sensor queda ciego
        (error del \SI{100}{\percent}), lo que obliga a incluir el
        \textbf{arranque ciego a lazo abierto} y la transición a lazo cerrado.
  \item \textbf{Retardo dependiente de la velocidad ($\tau \propto 1/n$):}
        aun por encima de $n_{\text{ciego}}$, el muestreo por eventos de cruce
        por cero y el promedio móvil degradan el margen de fase al bajar la
        velocidad.
  \item \textbf{Desenganche del acople magnético:} modelado como un acople
        mecánico/cinemático no lineal con torque máximo $T_{\max}$ y ángulo
        crítico $\delta_{\text{crit}}$ ($\SI{90}{\degree}$ para $N_a=1$ par de
        polos); al superarlo (p.\,ej.\ al agregar fluido o subir viscosidad),
        la barra se desincroniza aunque el motor siga regulando su propia
        velocidad.
  \item \textbf{Saturación en tensión del accionamiento:} el ciclo de trabajo
        queda acotado a $d\in[d_{\min},1]$ ($[0{,}05,\,1]$ en lazo cerrado) sobre
        una tensión de bus fija $V_{dc}$.
\end{enumerate}

\paragraph{Flujo metodológico.} Cada capítulo conecta directamente con los cuatro
pasos exigidos para el análisis de control:
\begin{enumerate}
  \item plantear el \textbf{modelo no lineal} con los fenómenos dominantes;
  \item modelar el \textbf{arranque ciego} hasta entrar en la zona observable;
  \item \textbf{linealizar} la planta en un punto de operación $(\omega_0,\delta^\ast,d_0)$
        y analizar el controlador PI lineal/discreto;
  \item comparar el diseño lineal contra la \textbf{planta no lineal} para
        delimitar su \textbf{rango de validez} (ceguera, saturación y desacople).
\end{enumerate}

Donde corresponde se toman los valores y algoritmos reales del firmware del
proyecto (STM32F103, \texttt{startup.c}, \texttt{speed\_sensor.c},
\texttt{motor\_control.c}, \texttt{rpm\_pi\_controller.c}).

\section{Diagrama del sistema}

\begin{center}
\begin{tikzpicture}[
  node distance=0.6cm and 0.55cm,
  blk/.style={draw,rounded corners,align=center,minimum height=1.0cm,
              font=\footnotesize,inner sep=4pt},
  >={Stealth}]
  \node[blk] (pi) {PI discreto\\+ sat.\ $d\in[d_{\min},1]$\\(cap.~\ref{cap:pi_lin})};
  \node[blk,right=of pi] (sw) {Modo\\Arranque /\\Lazo cerrado\\(cap.~\ref{cap:bldc})};
  \node[blk,right=of sw] (mot) {Motor BLDC\\eq.\ DC / SPWM\\(cap.~\ref{cap:bldc})};
  \node[blk,right=of mot] (acp) {Acople no lineal\\$T_{\max},\,\delta_{\text{crit}}$\\(cap.~\ref{cap:mecanica})};
  \node[blk,below=1.15cm of acp] (ag) {Agitador + fluido\\+ fricción\\(cap.~\ref{cap:cargas})};
  \node[blk,below=1.15cm of mot] (sen) {Sensor virtual ZC\\ceguera $<300\text{--}400$ rpm\\+ retardo $1/n$ (cap.~\ref{cap:medicion})};
  \draw[->] (pi) -- node[above,font=\scriptsize]{$d_k$} (sw);
  \draw[->] (sw) -- node[above,font=\scriptsize]{$d$ / $A,f_e$} (mot);
  \draw[->] (mot) -- node[above,font=\scriptsize]{$\theta_1,\omega_1$} (acp);
  \draw[->] (acp) -- node[right,font=\scriptsize]{$T_c(\delta)$} (ag);
  \draw[->] (mot.south) -- node[right,font=\scriptsize]{$\theta_1,\omega_1$} (sen.north);
  \draw[->] (sen.west) -| node[pos=0.25,below,font=\scriptsize]{$\hat n_k$, estado ZC} (pi.south);
  \node[left=0.45cm of pi,font=\footnotesize] (ref) {$n^\ast$};
  \draw[->] (ref) -- (pi);
\end{tikzpicture}
\end{center}

\section{Convenciones y notación}

\begin{table}[h]
\centering
\small
\begin{tabular}{@{}lll@{}}
\toprule
Símbolo & Significado & Unidad \\
\midrule
$\theta_1,\ \omega_1$ & ángulo y velocidad mecánica del rotor (lado motor) & \si{rad}, \si{rad/s} \\
$\theta_2,\ \omega_2$ & ángulo y velocidad mecánica del agitador & \si{rad}, \si{rad/s} \\
$\delta = \theta_1 - \theta_2$ & desfase angular en el acople magnético & \si{rad} \\
$p$ & pares de polos del motor (firmware: $p=2$, verificar en HDD) & -- \\
$\theta_e = p\,\theta_1,\ \omega_e = p\,\omega_1$ & ángulo y velocidad eléctricos & \si{rad}, \si{rad/s} \\
$n$ & velocidad en rpm, $n = 30\,\omega_1/\pi$ & \si{rpm} \\
$R,\ L_s$ & resistencia e inductancia por fase equivalentes & \si{\ohm}, \si{H} \\
$k_e,\ k_t$ & cte.\ de BEMF de fase y de torque ($k_t = 2k_e$ en six-step) & \si{V.s/rad}, \si{N.m/A} \\
$V_{dc},\ d$ & tensión de bus y ciclo de trabajo ($d\in[0{,}05,\,1]$) & \si{V}, -- \\
$J_1,\ J_2$ & inercia lado motor (rotor+imán) y lado agitador & \si{kg.m^2} \\
$N_a,\ \delta_{\text{crit}}$ & pares de polos del acople ($N_a=1$) y ángulo crítico ($\pi/(2N_a)$) & --, \si{rad} \\
$T_e,\ T_c,\ T_L,\ T_f$ & torque eléctrico, de acople, de carga del fluido, de fricción & \si{N.m} \\
$T_{\max}$ & torque máximo de desenganche (\emph{pull-out}) del acople & \si{N.m} \\
$n_{\text{ciego}}$ & velocidad mínima observable por cruces por cero ($\sim 300\text{--}400$) & \si{rpm} \\
$f_t,\ \Delta t_c$ & frecuencia del timer TIM2 (\SI{180}{kHz}) y período del PI (\SI{2}{ms}) & \si{Hz}, \si{s} \\
\bottomrule
\end{tabular}
\caption{Notación principal del modelo mínimo.}
\end{table}

La relación entre ángulo eléctrico y mecánico, $\theta_e = p\,\theta_1$, nace de
que un motor con $p$ pares de polos repite su distribución de campo $p$ veces
por vuelta mecánica. En el firmware actual $p=2$ por defecto
(\texttt{ESC\_DEFAULT\_POLE\_PAIRS}); como el prototipo usa un motor de disco
duro recuperado, este parámetro debe verificarse empíricamente sobre la maqueta
antes de cerrar los números finales.

% =====================================================================
\chapter{Modelo eléctrico del motor y arranque a lazo abierto}\label{cap:bldc}
% =====================================================================

\begin{idea}
No hace falta modelar en detalle la conmutación de los seis transistores ni el
\emph{cogging} para estudiar el control de velocidad y el desacople de la
barra. El motor cumple dos roles según la etapa de operación:
\begin{itemize}
  \item \textbf{Durante el arranque ciego:} el ESC impone un campo magnético
        rotante senoidal a lazo abierto (sin mirar el sensor) que ``engancha'' y
        arrastra al rotor desde el reposo hasta superar la zona ciega.
  \item \textbf{Una vez enganchado (six-step y lazo cerrado):} el motor se
        comporta como un motor de corriente continua equivalente con resistencia
        $2R$, inductancia $2L_s$ y constante $2k_e$.
\end{itemize}
\end{idea}

\section{Modelo equivalente en operación six-step}
Cuando el motor ya conmuta en seis pasos (\texttt{RUNNING} y
\texttt{CLOSEDLOOP}), conducen dos fases en serie en cada sector de
\SI{60}{\degree} eléctricos y la tercera queda flotante para medir BEMF. El
inversor (a $f_{PWM}=\SI{18}{kHz}$, mucho más rápida que la dinámica mecánica y
eléctrica) aplica una tensión media $d\,V_{dc}$ entre las dos fases activas. Por
Kirchhoff y Faraday:
\begin{equation}
  \boxed{\;d\,V_{dc} = 2R\,I + 2L_s\frac{\dd I}{\dd t} + 2k_e\,\omega_1\;}
  \label{eq:dcequiv}
\end{equation}
y por balance de potencia electromecánica ($P_{em} = 2k_e\,\omega_1\,I = T_e\,\omega_1$):
\begin{equation}
  \boxed{\;T_e = 2k_e\,I \equiv k_t\,I,\qquad k_t = 2k_e\;}
  \label{eq:torque_dc}
\end{equation}
donde $R$ y $L_s=L-M$ son la resistencia e inductancia síncrona por fase (que
se miden directamente entre terminales como $R_{LL}=2R$ y $L_{LL}=2L_s$).

\begin{lectura}
La tensión aplicada $d\,V_{dc}$ se reparte en tres términos: caída en el cobre
($2R\,I$), cambio de corriente en la inductancia ($2L_s\,\dot I$) y BEMF
($2k_e\,\omega_1$). En vacío y régimen permanente ($I\approx 0$), la velocidad
queda fijada por $\omega_1 \approx d\,V_{dc}/(2k_e)$. Cuando aparece torque de
carga, hace falta corriente, sube la caída $2RI$ y la velocidad cae con
pendiente $R/(2k_e^2)$: el controlador PI debe subir $d$ para sostenerla.
\end{lectura}

\paragraph{Saturación de tensión del inversor.} El ciclo de trabajo está
limitado físicamente por el bus disponible:
\begin{equation}
  d \in [d_{\min},\, 1], \qquad \text{con } d_{\min} = \frac{100}{2000} = 0{,}05
  \text{ en lazo cerrado.}
\end{equation}
En particular, $d\le 1$ impone una cota dura sobre la corriente y velocidad
máxima alcanzable: a tensión de bus $V_{dc}$ fija, no es posible pedir
cualquier torque ni cualquier aceleración. (El modelo trifásico completo
$v_{kn}$ y el \emph{cogging} se conservan como referencia en el
Apéndice~\ref{ap:motor_extendido}).

\section{Arranque ciego a lazo abierto (\texorpdfstring{\texttt{FOC\_STARTUP} $\to$ \texttt{RUNNING}}{FOC\_STARTUP -> RUNNING})}
\label{sec:arranque}

\begin{idea}
Por debajo de $\sim 300\text{--}\SI{400}{rpm}$ el detector de cruces por cero
está ciego. Por eso el firmware arranca ``a ciegas'': genera tres tensiones
senoidales lentas que crean un campo magnético giratorio cuya frecuencia y
amplitud crecen en rampa. Si la rampa es suficientemente suave, el rotor sigue
al campo (como un motor síncrono) hasta ganar velocidad y permitir que el
sensor empiece a ver la BEMF.
\end{idea}

En el firmware (\texttt{startup.c} y \texttt{hard\_config.h}), al ordenar el
arranque (\texttt{FOC\_STARTUP}) se ejecuta una rampa lineal de duración
$T_{\text{arr}}=\SI{3}{s}$:
\begin{align}
  f_e(t) &= f_{e,0} + (f_{e,f}-f_{e,0})\,\frac{t}{T_{\text{arr}}}, &
  f_{e,0}&=\SI{2,09}{Hz},\quad f_{e,f}=\SI{9,28}{Hz}, \label{eq:rampa_f}\\
  A(t) &= A_0 + (A_f-A_0)\,\frac{t}{T_{\text{arr}}}, &
  A_0&=0{,}20,\quad A_f=1{,}00, \label{eq:rampa_a}\\
  \theta_s(t) &= 2\pi\int_0^t f_e(\tau)\,\dd\tau.
\end{align}
Con modulación SPWM unipolar en las tres fases, la amplitud fundamental de
tensión fase-neutro aplicada es $V_{ph}(t) = \tfrac{1}{2}A(t)\,V_{dc}$. A las
bajas frecuencias eléctricas del arranque ($f_e < \SI{10}{Hz}$, donde
$2\pi f_e L_s \ll R$), la dinámica de la inductancia es despreciable y el torque
electromagnético de enganche síncrono puede aproximarse por:
\begin{equation}
  \boxed{\;T_{e,\text{arr}}(t) \approx \frac{3\,k_e}{2R}
  \Big(V_{ph}(t)\,\sin\!\big(\theta_s(t)-p\,\theta_1\big) - k_e\,\omega_1\Big)\;}
  \label{eq:torque_arranque}
\end{equation}
donde $\theta_s(t)-p\,\theta_1$ es el ángulo de carga eléctrico entre el campo
estatórico y el rotor.

\begin{lectura}
La ecuación~\eqref{eq:torque_arranque} explica en una línea cuándo funciona y
cuándo falla el arranque ciego:
\begin{itemize}
  \item Mientras el rotor acompaña al campo ($|\theta_s - p\theta_1|<\pi/2$), el
        seno genera un torque restaurador que acelera a $\omega_1$ hasta
        $\omega_{\text{campo}}(t)=2\pi f_e(t)/p$.
  \item Si la rampa $f_e(t)$ acelera demasiado rápido o la carga mecánica
        supera al par disponible $\frac{3k_e V_{ph}}{2R}$, el ángulo eléctrico
        supera $\pi/2$ y el motor \textbf{pierde sincronismo en el arranque}.
\end{itemize}
\end{lectura}

\paragraph{Transición al lazo cerrado (\emph{handoff}).} Al cumplirse
$t = T_{\text{arr}}$, \texttt{startup.c} detiene la senoidal, aplica un ciclo
de trabajo fijo $d_{\text{handoff}} = 0{,}45$ (\SI{45}{\percent}) en six-step y
pasa al estado \texttt{RUNNING} (ecuación~\eqref{eq:dcequiv}). Cuando la fase~W
acumula 4 cruces por cero consistentes, la máquina de estados pasa a
\texttt{READY} $\to$ \texttt{CLOSEDLOOP} y habilita el PI con inicialización sin
salto (\emph{bumpless}). Si transcurren \SI{4,5}{s} en \texttt{RUNNING} sin
enganchar cruces válidos (o si se dispara \emph{stall}), corta PWM \SI{0,5}{s}
y reintenta el arranque.

\begin{nota}
Con los valores por defecto del firmware ($f_{e,f}=\SI{9,28}{Hz}$) y $p=2$, la
velocidad síncrona al final de la rampa SPWM es
$n_{\text{sync}}(T_{\text{arr}}) = 60 f_{e,f}/p \approx \SI{278}{rpm}$, todavía
dentro o en el borde de la zona ciega ($300\text{--}\SI{400}{rpm}$). Es el
escalón a $d=0{,}45$ en \texttt{RUNNING} el que termina de acelerar el rotor
por encima de $n_{\text{ciego}}$ para enganchar la conmutación sensorless.
\end{nota}

% =====================================================================
\chapter{Mecánica de dos inercias y acople magnético no lineal}\label{cap:mecanica}
% =====================================================================

\begin{idea}
Hay dos cuerpos que giran: abajo, el rotor del motor con su imán impulsor
(inercia $J_1$); arriba, dentro del vaso, la cápsula o píldora agitadora
(inercia $J_2$). No hay un eje rígido entre ambos. Desde el punto de vista
mecánico-cinemático, el acople funciona como un resorte rotacional no lineal
que se estira al aparecer un desfase angular $\delta = \theta_1 - \theta_2$:
cuanto más se atrasa la cápsula, mayor es el torque que la arrastra,
\textbf{hasta llegar a un ángulo crítico $\delta_{\text{crit}}$}. Si se pasa de
ese ángulo, los imanes se sueltan, el torque cae de golpe y la barra patina.
\end{idea}

\section{Ecuaciones de Newton rotacionales}
Aplicando $\sum T = J\,\dot\omega$ a cada cuerpo por separado y llamando
$\delta = \theta_1 - \theta_2$ al desfase mecánico:
\begin{align}
  J_1\,\dot\omega_1 &= T_e - T_c(\delta) - T_{f1}(\omega_1), \label{eq:j1}\\
  J_2\,\dot\omega_2 &= T_c(\delta) - T_L(\omega_2) - T_{f2}(\omega_2), \label{eq:j2}\\
  \dot\delta &= \omega_1 - \omega_2 . \label{eq:delta}
\end{align}
Aquí $T_c(\delta)$ es el torque transmitido por el acople, $T_{f1}, T_{f2}$ son
las fricciones mecánicas (sección~\ref{sec:friccion}) y $T_L$ es la carga del
fluido (sección~\ref{sec:fluido}).

\section{Modelo cinemático/mecánico del acople magnético}\label{cap:acople}
Prescindiendo del detalle magnetostático (recogido en el
Apéndice~\ref{ap:acople_dipolo}), un acople entre imanes de $N_a$ pares de polos
queda caracterizado por dos magnitudes geométricas/empíricas claras:
\begin{itemize}
  \item el \textbf{torque máximo transmitible} (\emph{pull-out}) $T_{\max}$,
        que depende de los imanes y fuertemente de la distancia vertical entre
        el motor y el vaso;
  \item el \textbf{ángulo crítico de desenganche}
        \begin{equation}
          \delta_{\text{crit}} = \frac{\pi}{2 N_a},
        \end{equation}
        donde se alcanza el torque máximo ($N_a=1$ para una píldora magnética
        dipolar clásica $\Rightarrow \delta_{\text{crit}}=\SI{90}{\degree}$;
        $N_a=2$ si hubiera 4 polos $\Rightarrow \delta_{\text{crit}}=\SI{45}{\degree}$).
\end{itemize}
La ley periódica más simple que crece monótonamente desde $0$ hasta $T_{\max}$
en $\delta_{\text{crit}}$, cae al superarlo y se repite cada $2\pi/N_a$ es:
\begin{equation}
  \boxed{\;T_c(\delta) = T_{\max}\,\sin\!\left(\frac{\pi}{2}\,\frac{\delta}{\delta_{\text{crit}}}\right)
                       = T_{\max}\,\sin(N_a\,\delta).\;}
  \label{eq:acople}
\end{equation}
De forma conservadora, no se incluye amortiguamiento interno $c\,\dot\delta$ por
corrientes parásitas en el modelo principal, al no disponerse de medición para
ese coeficiente.

\begin{lectura}
Para $|\delta| < \delta_{\text{crit}}$, el desfase es un ``atraso'' elástico: si
la carga frena a la píldora, $\delta$ crece y el acople transmite más torque
para sostenerla. Para $|\delta| > \delta_{\text{crit}}$, el torque
\emph{disminuye} con el ángulo y al pasar $2\delta_{\text{crit}}$ ($\pi/N_a$)
cambia de signo: los polos opuestos se desencuentran y el acople pierde su
capacidad de arrastre.
\end{lectura}

\section{Estabilidad estática y desacople dinámico}
En régimen permanente a velocidad $\omega_0$, $\omega_1=\omega_2=\omega_0$ y
$\dot\delta=0$. De~\eqref{eq:j2}, el acople debe equilibrar la carga total
$T_{L,\text{tot}}(\omega_0) = T_L(\omega_0) + T_{f2}(\omega_0)$:
\begin{equation}
  T_{\max}\sin(N_a\delta^\ast) = T_{L,\text{tot}}(\omega_0).
  \label{eq:eq_acople}
\end{equation}
\begin{itemize}
  \item \textbf{Si $T_{L,\text{tot}} < T_{\max}$:} existen dos soluciones por
        período, pero perturbando $\delta=\delta^\ast+\tilde\delta$, el torque
        restaurador linealizado es $K_{\text{eff}}\,\tilde\delta$ con
        $K_{\text{eff}} = N_a T_{\max}\cos(N_a\delta^\ast)$. Sólo es estable la
        rama con $\cos(N_a\delta^\ast)>0$, es decir:
        \begin{equation}
          |\delta^\ast| < \delta_{\text{crit}} = \frac{\pi}{2N_a}.
        \end{equation}
  \item \textbf{Si $T_{L,\text{tot}} > T_{\max}$:} no existe equilibrio.
        $\delta(t)$ crece sin cota, el valor medio temporal
        $\langle T_{\max}\sin(N_a\delta)\rangle \to 0$ y el agitador
        \textbf{desacopla y se frena} ($\omega_2 \to 0$) aunque el motor siga
        girando a la consigna.
\end{itemize}

\paragraph{Desacople dinámico ante aceleraciones bruscas.} Aun con
$T_{L,\text{tot}} < T_{\max}$, un escalón rápido del motor puede provocar el
desenganche. Suponiendo $\dot\omega_1 \approx 0$ tras un cambio brusco de
velocidad, de~\eqref{eq:j2} y~\eqref{eq:delta}:
\begin{equation}
  J_2\,\ddot\delta = -T_{\max}\sin(N_a\delta) + T_{L,\text{tot}}.
  \label{eq:pendulo}
\end{equation}
Es la ecuación del péndulo con torque constante, con integral de energía:
\begin{equation}
  E = \tfrac{1}{2}J_2\dot\delta^2 - \frac{T_{\max}}{N_a}\cos(N_a\delta) - T_{L,\text{tot}}\,\delta = \text{cte}.
\end{equation}
El desenganche ocurre si la velocidad relativa inicial $\dot\delta$ aporta
energía cinética suficiente para cruzar la barrera potencial situada en el
equilibrio inestable $\delta_u = (\pi - \arcsin(T_{L,\text{tot}}/T_{\max}))/N_a$
(criterio de áreas iguales). Por eso una rampa suave de velocidad evita el
desacople frente a un escalón agresivo.

\section{Rigidez linealizada y modo torsional en el punto de operación}
A diferencia de linealizar en $\delta=0$ (válido sólo en vacío), alrededor de un
equilibrio con carga $\delta^\ast\ne 0$ la pendiente local de~\eqref{eq:acople}
es:
\begin{equation}
  \boxed{\;K_{\text{eff}}(\delta^\ast)
    = \left.\frac{\partial T_c}{\partial \delta}\right|_{\delta^\ast}
    = N_a\,T_{\max}\,\cos(N_a\,\delta^\ast)
    = N_a\sqrt{T_{\max}^2 - T_{L,\text{tot}}^2}\;.}
  \label{eq:keff}
\end{equation}
Despreciando temporalmente las fricciones y torques externos para mirar el
movimiento relativo $\ddot\delta = \dot\omega_1 - \dot\omega_2$, resulta el
oscilador torsional con frecuencia natural dependiente de la carga:
\begin{equation}
  \omega_n(\delta^\ast) = \sqrt{\frac{K_{\text{eff}}(\delta^\ast)}{J_{eq}}},
  \qquad J_{eq} = \frac{J_1 J_2}{J_1 + J_2}.
\end{equation}
\begin{lectura}
Cuanto más cerca está la carga de $T_{\max}$ ($\delta^\ast \to \delta_{\text{crit}}$),
más cae $K_{\text{eff}}$: el ``resorte magnético'' se ablanda hasta anularse en
el límite de desenganche. Ésta es una de las razones por las cuales el modelo
linealizado en vacío pierde validez al acercarse a carga máxima.
\end{lectura}

% =====================================================================
\chapter{Fricción y carga del fluido}\label{cap:cargas}
% =====================================================================

\section{Fricción mecánica: Coulomb más viscosa}\label{sec:friccion}
Para el modelo principal se adopta la descomposición clásica en fricción seca
constante (Coulomb) y fricción viscosa proporcional a la velocidad:
\begin{equation}
  \boxed{\;T_{f,i}(\omega_i) = T_{C,i}\tanh\!\left(\frac{\omega_i}{\omega_{0,f}}\right) + b_i\,\omega_i,
  \qquad i\in\{1,2\}\;.}
  \label{eq:friccion_simple}
\end{equation}
Se utiliza $\tanh(\omega_i/\omega_{0,f})$ con $\omega_{0,f}$ pequeño frente a
la velocidad de operación (p.\,ej.\ $\omega_{0,f}\approx \SI{5}{rad/s}$) en
lugar de $\sgn(\omega_i)$ únicamente para evitar una discontinuidad numérica en
$\omega_i=0$ durante la integración.

\paragraph{Por qué se omiten Stribeck y LuGre en el modelo principal.} Como el
sensor está ciego por debajo de $300\text{--}\SI{400}{rpm}$ y el motor arranca a
lazo abierto hasta superar el reposo, los parámetros adicionales de Stribeck o
LuGre no son identificables con la instrumentación disponible ni intervienen en
el régimen regulado; quedan documentados en el Apéndice~\ref{ap:friccion_extendida}.

\section{Ensayo de desaceleración libre (identificación)}
\begin{idea}
Si se acelera el motor en vacío (sin barra o en aire) y se corta la
alimentación ($d=0$, corriente nula), se detiene solo. Si la caída de
velocidad es lineal en el tiempo domina la fricción seca $T_C$; si es
exponencial domina la fricción viscosa $b$. Medir esa curva en telemetría
permite identificar la fricción del eje con datos reales.
\end{idea}
Con el PWM apagado ($T_e=0$) para $\omega>0$:
\begin{equation}
  J\dot\omega = -T_C - b\,\omega
  \quad\Longrightarrow\quad
  \omega(t) = \left(\omega_0 + \frac{T_C}{b}\right)e^{-t/\tau} - \frac{T_C}{b},
  \qquad \tau = \frac{J}{b},
  \label{eq:coastdown}
\end{equation}
donde $\omega_0$ es la velocidad al cortar el PWM. Ajustando la curva de
velocidad por encima de $n_{\text{ciego}}$ se estiman $T_C/J$ y $b/J$.

\section{Torque de carga del fluido}\label{sec:fluido}
La interacción de la píldora agitadora con el fluido dentro del vaso combina
dos efectos físicos de sentido claro:
\begin{itemize}
  \item un \textbf{término lineal (laminar/viscoso)} $b_{fl}\,\omega_2$,
        proporcional a la viscosidad $\mu$ del fluido;
  \item un \textbf{término cuadrático (inercial/turbulento)} $c_{fl}\,\omega_2|\omega_2|$,
        asociado a acelerar masa de fluido (proporcional a la densidad $\rho$ y
        al tamaño de la barra).
\end{itemize}
Sumando ambas contribuciones:
\begin{equation}
  \boxed{\;T_L(\omega_2) = b_{fl}\,\omega_2 + c_{fl}\,\omega_2|\omega_2|\;.}
  \label{eq:carga}
\end{equation}
\begin{nota}
A diferencia de tanques normalizados con deflectores (donde existen tablas de
número de potencia $N_p(\mathrm{Re})$ resumidas en el
Apéndice~\ref{ap:fluido_dimensional}), en un vaso de laboratorio sin deflectores
esos coeficientes publicados no aplican directamente. Para no introducir
constantes ajenas a la geometría real, $b_{fl}$ y $c_{fl}$ se tratan como
parámetros globales efectivos del medio (aire $\approx 0$, agua, mezclas
viscosas) a acotar o ajustar experimentalmente.
\end{nota}

\section{Escenario crítico: arranque en aire, agregado de fluido y desacople no observable}
\label{sec:aire_agua}
Para reproducir el ensayo clave discutido en la consulta del proyecto, se
parametriza la carga en función del tiempo o del medio:
\begin{equation}
  b_{fl}(t) = b_{fl,0} + \Delta b_{fl}\,\mathbb{1}(t - t_0), \qquad
  c_{fl}(t) = c_{fl,0} + \Delta c_{fl}\,\mathbb{1}(t - t_0).
\end{equation}
\begin{enumerate}
  \item Para $t < t_0$ el vaso está en vacío/aire ($b_{fl,0}\approx 0$,
        $c_{fl,0}\approx 0$): el sistema arranca a lazo abierto, engancha el
        lazo cerrado y motor y barra giran acoplados a $\omega_1=\omega_2=\omega_0$.
  \item En $t = t_0$ ingresa líquido o aumenta la viscosidad (agua / solución
        viscosa). Si la nueva carga supera $T_{\max}$ (o si el salto dinámico
        supera la barrera $\delta_u$ en~\eqref{eq:pendulo}), la píldora
        desacopla: $\delta$ crece, $\langle T_c\rangle \to 0$ y $\omega_2 \to 0$.
  \item \textbf{Consecuencia sobre la observabilidad:} al desaparecer el par
        transmitido $T_c$, la ecuación del motor~\eqref{eq:j1} queda
        desconectada del agitador ($J_1\dot\omega_1 \approx T_e - T_{f1}$). Como
        el firmware solo realimenta $\omega_1$ y no mide corriente ni movimiento
        de la barra, el controlador PI mantiene $\omega_1 \approx n^\ast$ (con
        apenas una perturbación transitoria leve) y \textbf{no detecta que el
        agitador se detuvo}.
\end{enumerate}

% =====================================================================
\chapter{Medición de velocidad sensorless}\label{cap:medicion}
% =====================================================================

\begin{idea}
El ESC \textbf{no mide la velocidad}: cronometra el tiempo entre cruces por cero
de la BEMF, que ocurren cada \SI{60}{\degree} eléctricos. A partir de esa
física aparecen las dos limitaciones que realmente condicionan el lazo en
nuestra maqueta:
\begin{enumerate}
  \item por debajo de $\sim 300\text{--}\SI{400}{rpm}$, la BEMF queda tapada por
        el ruido de conmutación y el umbral del comparador: \textbf{el sensor se
        queda ciego} (error del \SI{100}{\percent});
  \item por encima de ese umbral, la velocidad calculada siempre llega
        \textbf{atrasada}, porque mide el período que ya pasó y lo retiene
        hasta el cruce siguiente.
\end{enumerate}
\end{idea}

\section{Del período medido a la velocidad}
En conmutación six-step, cada fase queda flotante dos veces por ciclo eléctrico
y su BEMF cruza por cero en cada ventana. El tiempo entre dos cruces
consecutivos de la misma fase es medio ciclo eléctrico:
\begin{equation}
  T_{ZC} = \frac{1}{2f_e} = \frac{\pi}{p\,\omega_1} = \frac{30}{p\,n}\quad[\si{s}].
\end{equation}
El timer de captura TIM2 cuenta a $f_t = \SI{72}{MHz}/(399+1) = \SI{180}{kHz}$
(tick de \SI{5,56}{\micro s}), por lo que el período en ticks es
$P = f_t\,T_{ZC}$ y la velocidad estimada es:
\begin{equation}
  \boxed{\;n = \frac{30\, f_t}{p\, P}\;}
  \label{eq:period2rpm}
\end{equation}
que implementa directamente \texttt{period\_to\_rpm()} en
\texttt{speed\_sensor.c}.

\begin{ejemplo}
Con $f_t = \SI{180}{kHz}$ y $p=2$, resulta $n = 2{,}7\times10^6/P$:
\begin{center}
\begin{tabular}{@{}lrr@{}}
\toprule
Límite en firmware (\texttt{speed\_sensor.h} / \texttt{motor\_control.c}) & $P$ [ticks] & $n$ [rpm] \\
\midrule
Período máximo del timer de 16 bits & 65535 & 41 \\
\texttt{ZCP\_VALID\_MAX\_THRESHOLD} & 50000 & 54 \\
\texttt{SPEED\_MIN} (cota inferior del filtro) & 14000 & 193 \\
Límite empírico de ceguera en maqueta / \texttt{MIN\_RPM} & 9000--6750 & 300--400 \\
\texttt{SPEED\_MAX} (cota superior del filtro) & 200 & 13500 \\
\bottomrule
\end{tabular}
\end{center}
Si el motor físico de disco duro tiene $p > 2$, la relación período--rpm se
escala inversamente con $p$, pero el corte práctico impuesto por la amplitud de
BEMF frente al comparador sigue ocurriendo en el rango experimental de
$300\text{--}\SI{400}{rpm}$.
\end{ejemplo}

\section{Ceguera a baja velocidad (no linealidad dominante del sensor)}
Experimentalmente en la placa con comparador LM339 y el motor de disco duro,
por debajo de $n_{\text{ciego}} \approx 300\text{--}\SI{400}{rpm}$ no se detectan
cruces por cero utilizables:
\begin{itemize}
  \item \textbf{Error de medición del \SI{100}{\percent}:} no hay eventos de
        captura válidos; el lazo cerrado no dispone de realimentación de
        velocidad real.
  \item \textbf{Comportamiento en firmware:} durante un breve lapso la rutina
        \texttt{filtro\_\allowbreak media\_\allowbreak movil\_\allowbreak zc()}
        retiene el último período válido; pero si transcurren más de
        $\tau_{\text{stall}} = \SI{200}{ms}$
        (\texttt{TIMEOUT\_\allowbreak MOTOR\_\allowbreak STALL\_\allowbreak MS},
        verificado cada \SI{25}{ms} hasta 3~conteos) sin cruces en la fase~W,
        \texttt{check\_\allowbreak motor\_\allowbreak status()} marca
        \texttt{motor\_stalled = true}, corta el PWM y relanza el arranque a
        lazo abierto (\texttt{foc\_startup()}).
\end{itemize}
Por ello, en el modelo de simulación la condición
\begin{equation}
  n_1 < n_{\text{ciego}} \quad (\approx 300\text{--}\SI{400}{rpm})
\end{equation}
bloquea los nuevos eventos ZC y, tras $\tau_{\text{stall}}$, fuerza la señal de
pérdida de medición ($\hat n_k \to 0$ y bandera de \emph{stall}), evitando la
ficción de retener eternamente la última velocidad válida cuando el rotor se
frenó.

\section{Retardo por muestreo por eventos (\texorpdfstring{$\tau \propto 1/n$}{tau proporcional a 1/n})}\label{sec:retardo}
Aun por encima de $n_{\text{ciego}}$, el sensado introduce un retardo variable
con la velocidad debido a tres contribuciones acumulativas:
\begin{enumerate}
  \item \textbf{Promedio móvil de $M$ períodos:} el firmware promedia
        $M=\texttt{ZCP\_BUFFER\_SIZE}=2$ períodos, lo que centra la ventana con
        un retardo medio de $M T_{ZC}/2$.
  \item \textbf{Retención entre cruces (ZOH por eventos):} entre un cruce y el
        siguiente el dato queda congelado; retardo medio $T_{ZC}/2$.
  \item \textbf{Muestreo del PI} cada $\Delta t_c = \SI{2}{ms}$: retardo medio
        $\Delta t_c/2$.
\end{enumerate}
Sumando las tres contribuciones y sustituyendo $T_{ZC}=30/(p\,n)$:
\begin{equation}
  \boxed{\;\tau(n) \approx \frac{M+1}{2}\,T_{ZC} + \frac{\Delta t_c}{2}
         = \frac{15\,(M+1)}{p\,n} + \frac{\Delta t_c}{2}\;.}
  \label{eq:retardo}
\end{equation}

\begin{ejemplo}
Con $M=2$, $p=2$ y $\Delta t_c=\SI{2}{ms}$: a \SI{400}{rpm},
$\tau \approx \SI{57,3}{ms}$; a \SI{1000}{rpm}, $\tau \approx \SI{23,5}{ms}$; a
\SI{3000}{rpm}, $\tau \approx \SI{8,5}{ms}$. El retardo crece como $1/n$.
\end{ejemplo}

\paragraph{Consecuencia para el margen de fase.} Un retardo $e^{-s\tau}$ aporta
fase $-\omega\tau$ sin alterar el módulo. Si el lazo diseñado a velocidad nominal
tiene frecuencia de cruce $\omega_c$ y un margen $\Delta\phi$ disponible antes
de volverse oscilatorio, debe cumplirse $\omega_c\,\tau(n) < \Delta\phi$:
\begin{equation}
  \boxed{\;n_{\min,\text{estab}} \approx \frac{15\,(M+1)\,\omega_c}{p\,\Delta\phi}\;.}
\end{equation}
Así, un PI fijo bien ajustado a \SI{2000}{rpm} puede empezar a oscilar al bajar
la consigna antes incluso de tocar el límite de ceguera. (Los efectos
secundarios de alta velocidad —cuantización $\delta n\propto n^2$, retardo RC,
\emph{blanking} y \emph{jitter} fino— se resumen en el
Apéndice~\ref{ap:sensor_extendido}).

\section{Algoritmo del sensor virtual para Simulink}\label{sec:sensor_virtual}
El bloque recibe $(\theta_1, \omega_1, t)$ y reproduce los efectos dominantes
(ceguera $< n_{\text{ciego}}$, timeout de \emph{stall}, retardo por eventos,
aritmética entera del timer de 16 bits y filtro $M=2$):
\begin{small}
\begin{verbatim}
function [n_hat, stalled] = sensor_virtual(theta_m, omega_m, t, ...
                                           p, ft, M, n_blind, tau_stall)
% Sensor virtual ZC: ceguera a baja rpm + timeout + retardo por eventos.
persistent last_k last_stamp periods idx count n_out last_valid_t
if isempty(last_k)
    last_k = zeros(3,1);  last_stamp = -ones(3,1);
    periods = zeros(3,M); idx = ones(3,1);
    count = zeros(3,1);   n_out = 0;  last_valid_t = 0;
end
n_real = abs(omega_m)*30/pi;
theta_e = p*theta_m;
phi = [0, 2*pi/3, -2*pi/3];
for k = 1:3
    kk = floor((theta_e - phi(k))/pi);
    if kk == last_k(k), continue; end
    last_k(k) = kk;
    if n_real < n_blind                  % zona ciega: no detecta cruces
        continue;
    end
    stamp = mod(round(t*ft), 65536);     % captura timer de 16 bits
    if last_stamp(k) >= 0
        P = mod(stamp - last_stamp(k), 65536);
        if P > 50 && P < 50000
            periods(k,idx(k)) = P;
            idx(k) = mod(idx(k), M) + 1;
            count(k) = min(count(k)+1, M);
            if count(k) == M
                Pavg = floor(mean(periods(k,:)));
                ok = all(abs(periods(k,:) - Pavg) <= 0.25*Pavg);
                if ok && Pavg >= 200 && Pavg <= 14000
                    n_out = floor(30*ft/(p*Pavg));
                    last_valid_t = t;
                end
            end
        end
    end
    last_stamp(k) = stamp;
end
if (t - last_valid_t) > tau_stall        % perdida de senal (>200 ms)
    n_out = 0;
    stalled = 1;
else
    stalled = 0;
end
n_hat = n_out;
\end{verbatim}
\end{small}

% =====================================================================
\chapter{Controlador PI discreto, linealización y rango de validez}\label{cap:pi_lin}
% =====================================================================

\section{Controlador PI discreto del firmware}
El controlador opera cada $\Delta t_c = \SI{2}{ms}$ sobre el error en rpm
$e_k = n^\ast - \hat n_k$, generando una salida en cuentas canónicas
$u_k \in [u_{\min}, u_{\max}] = [100,\,2000]$ (que en la planta corresponde al
ciclo de trabajo $d_k = u_k/2000 \in [0{,}05,\,1]$).

\paragraph{Discretización trapezoidal (Tustin).} A partir del PI continuo
$u(t) = K_p e(t) + K_i\int_0^t e(\tau)\,\dd\tau$, aproximando el integrador
por trapecios de paso $\Delta t_c$:
\begin{equation}
  I_k = I_{k-1} + K_i\,\Delta t_c\,\frac{e_k + e_{k-1}}{2}, \qquad
  u_k^{\text{lin}} = K_p\,e_k + I_k .
  \label{eq:pi_tustin}
\end{equation}
En transformada $\mathcal Z$:
\begin{equation}
  \boxed{\;C(z) = \frac{U(z)}{E(z)} = K_p + \frac{K_i\,\Delta t_c}{2}\,\frac{z+1}{z-1}\;.}
  \label{eq:cz}
\end{equation}

\paragraph{Saturación y anti-windup por \emph{clamping}.} Para evitar que el
integrador siga acumulando error cuando el inversor alcanza el límite
$[u_{\min},u_{\max}]$ (por ejemplo durante un bloqueo o un escalón grande), el
firmware recorta el estado integral antes de saturar la salida:
\begin{equation}
\begin{aligned}
  I_k &\leftarrow \clamp\!\left(I_k,\; u_{\min} - K_p e_k,\; u_{\max} - K_p e_k\right), \\
  u_k &= \clamp\!\left(K_p e_k + I_k,\; u_{\min},\; u_{\max}\right),
  \qquad d_k = \frac{u_k}{2000}.
\end{aligned}
\label{eq:antiwindup}
\end{equation}
(El detalle de implementación entera en punto fijo Q16.16/Q30 del STM32 se
documenta en el Apéndice~\ref{ap:q16}).

\section{Modelo no lineal completo y equilibrio de operación}
En operación normal tras el arranque (modo \texttt{CLOSEDLOOP}), la dinámica
completa de la planta reúne las ecuaciones continuas~\eqref{eq:dcequiv},
\eqref{eq:torque_dc}, \eqref{eq:j1}--\eqref{eq:delta}, \eqref{eq:acople},
\eqref{eq:friccion_simple} y~\eqref{eq:carga}:
\begin{align}
  2L_s\,\dot I &= d_k V_{dc} - 2R\,I - 2k_e\,\omega_1, \\
  J_1\,\dot\omega_1 &= 2k_e I - T_{\max}\sin(N_a\delta) - T_{C1}\tanh(\omega_1/\omega_{0,f}) - b_1\omega_1, \\
  J_2\,\dot\omega_2 &= T_{\max}\sin(N_a\delta) - b_{fl}\omega_2 - c_{fl}\omega_2|\omega_2| - T_{C2}\tanh(\omega_2/\omega_{0,f}) - b_2\omega_2, \\
  \dot\delta &= \omega_1 - \omega_2, \qquad \dot\theta_1 = \omega_1 .
\end{align}

\paragraph{Punto de equilibrio $(\omega_0, I_0, \delta^\ast, d_0)$.} Fijada una
velocidad de operación $\omega_0 > \omega_{0,f}$ (donde $\tanh\approx 1$,
$n_0 = \frac{30}{\pi}\omega_0 > n_{\text{ciego}}$) y anulando derivadas:
\begin{align}
  T_{L,\text{tot}}(\omega_0) &= b_{fl}\omega_0 + c_{fl}\omega_0|\omega_0| + T_{C2} + b_2\omega_0, \\
  \delta^\ast &= \frac{1}{N_a}\arcsin\!\left(\frac{T_{L,\text{tot}}(\omega_0)}{T_{\max}}\right)
  \qquad\big(\text{existe sólo si } T_{L,\text{tot}}(\omega_0) < T_{\max}\big), \\
  I_0 &= \frac{T_{L,\text{tot}}(\omega_0) + T_{C1} + b_1\omega_0}{2k_e}, \\
  d_0 &= \frac{2R\,I_0 + 2k_e\,\omega_0}{V_{dc}}
  \qquad\big(\text{requiere } d_0 \in [d_{\min}, 1]\big).
\end{align}

\section{Linealización alrededor de \texorpdfstring{$(\omega_0,\delta^\ast,d_0)$}{(omega0, delta*, d0)}}
Definiendo variables incrementales $\tilde x = x - x_0$ para $\omega_1,\omega_2 > 0$
(en torno a $\omega_0 \gg \omega_{0,f}$ la fricción seca $T_{C,i}$ es constante y
su derivada es nula):
\begin{align}
  B_{1,\text{eq}} &= b_1, \\
  B_{2,\text{eq}}(\omega_0) &= b_2 + b_{fl} + 2\,c_{fl}\,\omega_0, \\
  K_{\text{eff}}(\delta^\ast) &= N_a\,T_{\max}\cos(N_a\,\delta^\ast)
                               = N_a\sqrt{T_{\max}^2 - T_{L,\text{tot}}(\omega_0)^2}.
\end{align}
El sistema linealizado en espacio de estados $\dot{\tilde x} = A\,\tilde x + B\,\tilde d$,
$\tilde n_1 = C\,\tilde x$ (con $\tilde x = [\tilde I,\ \tilde\omega_1,\ \tilde\omega_2,\ \tilde\delta]^\mathsf{T}$
y salida en rpm) es:
\begin{equation}
  A = \begin{bmatrix}
    -\dfrac{R}{L_s} & -\dfrac{k_e}{L_s} & 0 & 0 \\[6pt]
    \dfrac{2k_e}{J_1} & -\dfrac{B_{1,\text{eq}}}{J_1} & 0 & -\dfrac{K_{\text{eff}}}{J_1} \\[6pt]
    0 & 0 & -\dfrac{B_{2,\text{eq}}}{J_2} & \dfrac{K_{\text{eff}}}{J_2} \\[6pt]
    0 & 1 & -1 & 0
  \end{bmatrix},
  \qquad
  B = \begin{bmatrix}
    \dfrac{V_{dc}}{2L_s} \\[6pt] 0 \\[4pt] 0 \\[4pt] 0
  \end{bmatrix},
  \qquad
  C = \begin{bmatrix} 0 & \dfrac{30}{\pi} & 0 & 0 \end{bmatrix}.
  \label{eq:ss_lineal}
\end{equation}

\paragraph{Reducción para análisis a mano (lápiz y papel).} Como la constante
eléctrica $\tau_e = L_s/R$ suele ser mucho menor que las mecánicas, puede
tomarse corriente cuasi-estática $\tilde I \approx (V_{dc}\tilde d - 2k_e\tilde\omega_1)/(2R)$.
Eso introduce el amortiguamiento eléctrico $B_e = 2k_e^2/R$ sobre el rotor,
llamando $B_1' = B_{1,\text{eq}} + 2k_e^2/R$, y da la función de transferencia
de 2 inercias:
\begin{equation}
  G_2(s) = \frac{\Omega_1(s)}{D(s)}
  = \frac{\dfrac{k_e V_{dc}}{R}\,(J_2 s^2 + B_{2,\text{eq}} s + K_{\text{eff}})}
         {(J_1 s + B_1')(J_2 s^2 + B_{2,\text{eq}} s + K_{\text{eff}}) + K_{\text{eff}}(J_2 s + B_{2,\text{eq}})}.
  \label{eq:g2s}
\end{equation}
\begin{lectura}
La ecuación~\eqref{eq:g2s} muestra dos casos límite inmediatos a mano:
\begin{itemize}
  \item \textbf{Acople rígido ($K_{\text{eff}}\to\infty$, operación lejos del
        límite de desenganche):} el sistema colapsa a un primer orden con
        $J_{\text{tot}} = J_1 + J_2$ y $B_{\text{tot}} = B_{1,\text{eq}} + B_{2,\text{eq}}(\omega_0)$,
        ideal para ajustar el PI a mano:
        \begin{equation}
          G_1(s) = \frac{\Omega_1(s)}{D(s)} = \frac{K_{dc}}{\tau_{\text{mec}} s + 1},
          \qquad
          K_{dc} = \frac{V_{dc}/(2k_e)}{1 + \dfrac{R\,B_{\text{tot}}}{2k_e^2}},
          \qquad
          \tau_{\text{mec}} = \frac{J_{\text{tot}}}{B_{\text{tot}} + \dfrac{2k_e^2}{R}}.
          \label{eq:g1s}
        \end{equation}
        (Recordar multiplicar por $\frac{30}{\pi}\frac{1}{2000}$ si se conecta
        directamente al PI que mide error en rpm y entrega cuentas sobre 2000.)
  \item \textbf{Acople perdido ($K_{\text{eff}}\to 0$, en $\delta^\ast\to\delta_{\text{crit}}$ o tras desenganche):}
        el par de polos y ceros del agitador $(J_2 s^2 + B_{2,\text{eq}} s)$ se
        cancela en~\eqref{eq:g2s}, dejando únicamente la planta del motor solo
        $\frac{k_e V_{dc}/R}{J_1 s + B_1'}$: la dinámica de la barra desaparece
        por completo de la salida $\omega_1$!
\end{itemize}
\end{lectura}

Incluyendo el retardo de medición~\eqref{eq:retardo} en el punto de operación,
el lazo linealizado completo a analizar es $L(s) = \frac{30}{2000\pi}\,C(s)\,G_2(s)\,e^{-s\tau(n_0)}$
(o en tiempo discreto mediante la discretización de $G(s)\,e^{-s\tau(n_0)}$ con
$C(z)$).

\section{Rango de validez de la aproximación lineal}
El controlador lineal diseñado con~\eqref{eq:g1s}--\eqref{eq:g2s} coincide con la
planta real sólo dentro de una ventana de operación acotada. En la planta no
lineal deja de funcionar por cuatro fronteras concretas:

\begin{table}[h]
\centering
\small
\begin{tabular}{@{}>{\raggedright\arraybackslash}p{3.2cm}>{\raggedright\arraybackslash}p{4.3cm}>{\raggedright\arraybackslash}p{7.3cm}@{}}
\toprule
Frontera no lineal & Condición de ruptura & Qué ocurre en la planta no lineal \\
\midrule
Zona ciega del sensor & $n_1 < n_{\text{ciego}}$ ($\sim 300\text{--}\SI{400}{rpm}$) & Error \SI{100}{\percent}, timeout \SI{200}{ms} $\to$ \emph{stall} y reinicio \\
Retardo variable $1/n$ & $n_1 < n_{\min,\text{estab}}$ & Cae el margen de fase; ciclo oscilatorio sin cambiar el PI \\
Saturación de tensión & $d_k$ choca en $0{,}05$ ó $1{,}0$ & Lazo abierto temporal; el anti-windup frena el integrador \\
Desacople del acople & $T_{L,\text{tot}} > T_{\max}$ ó $\delta > \delta_u$ & $K_{\text{eff}}\to 0$, $\omega_2\to 0$; el PI sigue en $\omega_1\approx n^\ast$ sin detectarlo \\
\bottomrule
\end{tabular}
\caption{Fronteras que definen el rango de validez del modelo lineal.}
\end{table}

% =====================================================================
\chapter{Parámetros, implementación en Simulink y plan de pruebas}\label{cap:escalar}
% =====================================================================

\section{Origen real de cada parámetro}
Para evitar parámetros inventados sin sustento experimental, cada magnitud del
modelo mínimo tiene asignado su método de obtención o acotación conservadora:

\begin{table}[h]
\centering
\small
\begin{tabular}{@{}ll>{\raggedright\arraybackslash}p{6.8cm}@{}}
\toprule
Parámetro & Origen & Cómo se obtiene / justifica \\
\midrule
$f_t,\Delta t_c,f_{PWM},M,d_{\min}$ & firmware & $\SI{180}{kHz}$, $\SI{2}{ms}$, $\SI{18}{kHz}$, $2$, $0{,}05$ (exactos del código) \\
$T_{\text{arr}},f_{e,0},f_{e,f},A_0,A_f$ & firmware & $\SI{3}{s}$, $\SI{2,09}{Hz}$, $\SI{9,28}{Hz}$, $0{,}20$, $1{,}00$ (\texttt{hard\_config.h}) \\
$K_p,\ K_i$ & firmware & Ganancias del PI discreto (ej.\ $K_p=0{,}28$, $K_i=1{,}00$) \\
$R,\ L_s$ & medición & Medidos con multímetro y osciloscopio sobre las fases del HDD \\
$k_e,\ p$ & medición / inspección & $k_e$ midiendo BEMF con giro externo en osciloscopio; verificar $p$ del HDD ($p=2$ default) \\
$V_{dc}$ & fuente / medición & Tensión de alimentación medida en bornes del puente L298N \\
$J_1,\ J_2$ & CAD + balanza & Geometría de soporte, imán impulsor y píldora agitadora \\
$T_C,\ b$ & ensayo en vacío & Curva de desaceleración libre~\eqref{eq:coastdown} con telemetría \\
$N_a,\ \delta_{\text{crit}}$ & geometría & Pares de polos del imán ($N_a=1 \Rightarrow \delta_{\text{crit}}=\SI{90}{\degree}$) \\
$T_{\max}$ & medición estática / barrido & Traba mecánica + brazo/dinamómetro (o ángulo límite), o barrido conservador \\
$b_{fl},\ c_{fl}$ & ajuste empírico & Aire ($\approx 0$), agua y fluido viscoso; barrido para capturar frontera de desacople \\
$n_{\text{ciego}},\ \tau_{\text{stall}}$ & banco + firmware & $300\text{--}\SI{400}{rpm}$ observado en comparador; $\tau_{\text{stall}}=\SI{200}{ms}$ \\
\bottomrule
\end{tabular}
\caption{Trazabilidad de los parámetros del modelo mínimo.}
\end{table}

\section{Implementación escalar en Simulink}
\begin{idea}
Un diagrama de Simulink sin bloques ocultos (sin usar cajas negras de Simscape
que oculten los signos o parámetros internos) es una ecuación diferencial
dibujada: todo lo que aparece derivado sale de un \texttt{Integrator}, y cada
cable tiene significado físico comprobable a mano.
\end{idea}

\begin{table}[h]
\centering
\small
\begin{tabular}{@{}llll@{}}
\toprule
Estado & Integrador de & Significado físico & Cond.\ inicial \\
\midrule
$I$ & $\dot I$ (ec.~\eqref{eq:dcequiv}) & corriente de fase equivalente & $0$ \\
$\omega_1$ & $\dot\omega_1$ (ec.~\eqref{eq:j1}) & velocidad angular del motor & $0$ \\
$\omega_2$ & $\dot\omega_2$ (ec.~\eqref{eq:j2}) & velocidad angular de la píldora & $0$ \\
$\theta_1$ & $\omega_1$ & posición angular del motor (para sensor y arranque) & $0$ \\
$\delta$ & $\omega_1 - \omega_2$ & desfase en el acople (evita restar ángulos grandes) & $0$ \\
\bottomrule
\end{tabular}
\caption{Estados continuos del modelo reducido.}
\end{table}

\paragraph{Estructura de bloques.}
\begin{itemize}
  \item \textbf{Selector de modo (\texttt{FOC\_STARTUP} vs \texttt{CLOSEDLOOP}):}
        para $t < T_{\text{arr}}$, un bloque \texttt{MATLAB Function} genera
        $\theta_s(t)$, $A(t)$ y $T_{e,\text{arr}}$ según~\eqref{eq:torque_arranque};
        al finalizar la rampa, conmuta a la ecuación~\eqref{eq:dcequiv} con
        $d = 0{,}45$ hasta que el sensor virtual valide cruces consistentes
        (\texttt{stalled == 0} y $n_1 \ge n_{\text{ciego}}$) y habilita el PI.
  \item \textbf{Subsistema eléctrico (six-step):} \texttt{Product} ($d\cdot V_{dc}$),
        realimentación $-2R\,I$ y $-2k_e\,\omega_1$, \texttt{Gain} $1/(2L_s)$
        e \texttt{Integrator} $\to I$ (opcionalmente limitado a $I\ge 0$ si no
        hay regeneración), y \texttt{Gain} $2k_e \to T_e$.
  \item \textbf{Acople mecánico no lineal:} \texttt{Sum} $\omega_1-\omega_2$
        $\to$ \texttt{Integrator} $\to \delta$ $\to$ \texttt{Fcn}
        \verb|Tmax*sin(Na*u)| $\to T_c$ (resta en $J_1$, suma en $J_2$).
  \item \textbf{Fricción y fluido:} bloques \texttt{Fcn} sin estado propio con
        expresiones \verb|Tc*tanh(u/w0f) + b*u| (fricción) y
        \verb|bfl(t)*u + cfl(t)*u*abs(u)| (fluido).
  \item \textbf{Sensor virtual y PI discreto:} el sensor de la
        sección~\ref{sec:sensor_virtual} como \texttt{MATLAB Function}, y el PI
        discreto~\eqref{eq:pi_tustin}--\eqref{eq:antiwindup} con \emph{Sample time}
        $= \SI{0,002}{s}$.
  \item \textbf{Solver:} paso fijo \texttt{ode4} (RK4) con $h = 1/f_t \approx \SI{5,56}{\micro s}$,
        de modo que las capturas queden cuantizadas al tick real de TIM2 y el
        PI caiga exactamente cada 360 pasos.
\end{itemize}

\section{Plan mínimo de simulaciones de verificación}
\begin{enumerate}
  \item \textbf{Validación analítica en lazo abierto:} motor sin carga con $d$
        constante $\to$ verificar $\omega_1 \to d\,V_{dc}/(2k_e)$; luego corte
        $d=0$ $\to$ verificar desaceleración libre~\eqref{eq:coastdown}.
  \item \textbf{Arranque ciego a lazo abierto:} verificar que la rampa SPWM
        saca al sistema desde $\omega_1=\omega_2=0$ en aire, cruza la zona
        ciega ($300\text{--}\SI{400}{rpm}$) y engancha el PI sin desacoplar la
        píldora; luego probar una rampa más rápida para observar pérdida de
        sincronismo en el arranque o en $\delta$.
  \item \textbf{Lineal vs.\ no lineal (rango de validez):} comparar respuesta al
        escalón del modelo lineal~\eqref{eq:g1s}--\eqref{eq:g2s} contra el
        modelo no lineal en un punto nominal (p.\,ej.\ \SI{1000}{rpm}); luego
        reducir consigna hacia $300\text{--}\SI{400}{rpm}$ (oscilación por
        retardo + ceguera) y exigir escalones grandes (saturación de $d$ y
        desacople dinámico).
  \item \textbf{Experimento aire $\to$ líquido / cambio de viscosidad:}
        arrancar en vacío, estabilizar en lazo cerrado y aplicar en $t=t_0$ el
        escalón de $b_{fl}, c_{fl}$ hasta superar $T_{\max}$: verificar que
        $\omega_2 \to 0$ mientras el controlador sigue regulando
        $\omega_1 \approx n^\ast$ sin advertir el desacople.
\end{enumerate}

% =====================================================================
\appendix
\chapter{Material de referencia y extensiones (fuera del modelo principal)}
\label{ap:referencia}
% =====================================================================
\noindent\textit{Nota:} Las secciones de este apéndice recogen deducciones
extendidas y detalles de implementación que \textbf{no forman parte del modelo
mínimo ni se asumen como hipótesis activas de la simulación principal}, salvo
que un ensayo específico justifique incorporarlos.

\section{Modelo trifásico completo y torque de cogging}\label{ap:motor_extendido}
En el modelo detallado de Pillay y Krishnan~\cite{pillay1989}, cada fase
$k\in\{a,b,c\}$ de un motor en estrella sin neutro ($i_a+i_b+i_c=0$, $L_s=L-M$)
satisface:
\begin{equation}
  v_{kn} = R\,i_k + L_s\frac{\dd i_k}{\dd t} + k_e f_k(\theta_e)\,\omega_1,
  \qquad
  v_{nN} = \frac{1}{3}\sum_k v_{kN} - \frac{1}{3}\sum_k e_k,
  \qquad
  T_e = k_e\sum_k f_k(\theta_e)\,i_k .
\end{equation}
Para $f_a(\theta_e)$ trapezoidal unitaria con dos fases activas conduciendo
$\pm I$ en los planos ($\pm 1$), se recupera $T_e = 2k_e I$ y~\eqref{eq:dcequiv}.

El torque de \emph{cogging} (interacción magnética sin corriente entre imanes y
ranuras $Q$) es periódico en $\theta_1$ con $N_c=\operatorname{mcm}(Q,2p)$:
$T_{cog}(\theta_1) = \sum_{h\ge 1} T_{cog,h}\sin(hN_c\theta_1 + \varphi_h)$~\cite{hanselman2006}.
Se omite del modelo principal por carecer de medición de $T_{cog,h}$.

\section{Origen dipolar del acople y amortiguamiento parásito}\label{ap:acople_dipolo}
A primer orden, un imán de momento $\mathbf m$ en un campo $\mathbf B$ tiene
energía $U=-\mathbf m\cdot\mathbf B = -mB\cos\alpha$ y torque
$\tau = -\dd U/\dd\alpha = -mB\sin\alpha$~\cite{griffiths2017}. Generalizando
a un acople periódico de $N_a$ pares de polos por serie de Fourier par,
$U(\delta)\approx U_0 - U_1\cos(N_a\delta)$, por trabajos virtuales
$T_c = \partial U/\partial\delta = N_a U_1\sin(N_a\delta) = T_{\max}\sin(N_a\delta)$~\cite{montague2012,ravaud2009}.
En acoples con carcasas conductoras puede aparecer un término viscoso adicional
$c\,\dot\delta$ por corrientes parásitas, omitido por conservadurismo.

\section{Modelos extendidos de fricción (Stribeck, LuGre, Petroff)}\label{ap:friccion_extendida}
Para reproducir \emph{stick-slip} cerca de $\omega=0$, la formulación de
\emph{Rotational Friction} de Simscape~\cite{simscape-friction,armstrong1994}
agrega al modelo Coulomb + viscoso el pico de despegue $T_{brk} > T_C$:
\begin{equation}
  T_f(\omega) = \sqrt{2e}\,(T_{brk}-T_C)\,e^{-(\omega/\omega_{St})^2}\frac{\omega}{\omega_{St}}
              + T_C\tanh\!\left(\frac{\omega}{\omega_{Coul}}\right) + b\,\omega,
\end{equation}
con $\omega_{St}=\sqrt2\,\omega_{brk}$ y $\omega_{Coul}=\omega_{brk}/10$. El
modelo dinámico LuGre~\cite{canudas1995} agrega un estado interno de asperezas
$\dot z = \omega - \frac{\sigma_0|\omega|}{g(\omega)}z$,
$T_f=\sigma_0 z + \sigma_1\dot z + b\,\omega$. Por su parte, un buje lubricado
de radio $r$, largo $\ell$ y huelgo $h$ (ley de Petroff) explica la fricción
viscosa lineal como $b = 2\pi\mu r^3\ell/h$.

\section{Análisis dimensional del fluido (Buckingham y \texorpdfstring{$N_p$}{Np})}\label{ap:fluido_dimensional}
Para un agitador de diámetro $D$ girando a $N=\omega_2/(2\pi)$ rev/s en un
fluido de densidad $\rho$ y viscosidad $\mu$, sin formación de vórtice, el
teorema $\Pi$ de Buckingham~\cite{buckingham1914} reduce $T_L=\phi(N,D,\rho,\mu)$
al número de potencia $N_p = 2\pi T_L/(\rho N^2 D^5)$ en función de
$\mathrm{Re} = \rho N D^2/\mu$~\cite{rushton1950,paul2004,mccabe2005}:
\begin{itemize}
  \item régimen laminar ($\mathrm{Re}\lesssim 10$): $N_p\mathrm{Re}=K_L \implies
        T_L = \frac{K_L \mu D^3}{4\pi^2}\,\omega_2 = b_{fl}\omega_2$;
  \item régimen turbulento ($\mathrm{Re}\gtrsim 10^4$): $N_p\to N_{p,\infty} \implies
        T_L = \frac{N_{p,\infty}\rho D^5}{8\pi^3}\,\omega_2^2 = c_{fl}\omega_2^2$.
\end{itemize}
Como $K_L$ y $N_{p,\infty}$ tabulados corresponden a tanques con deflectores, en
el modelo principal se identifican o barren directamente $b_{fl}$ y $c_{fl}$.

\section{Imperfecciones secundarias del sensor ZC}\label{ap:sensor_extendido}
Además de la ceguera a baja velocidad y el retardo $\tau\propto 1/n$ del
Capítulo~\ref{cap:medicion}, el sensor presenta efectos de segundo orden:
\begin{itemize}
  \item \textbf{Cuantización de ticks (alta velocidad):} para $P$ entero,
        $\delta n \approx \frac{p\,n^2}{30\,f_t}$ (a \SI{1000}{rpm} con $p=2$,
        apenas $\SI{0,37}{rpm}$)~\cite{ohmae1982}.
  \item \textbf{\emph{Jitter} por ruido de tensión:} con ruido $\sigma_v$ en
        la entrada del comparador, la pendiente en el cruce es
        $\dot e = \kappa k_e p\,\omega_1^2$, dando $\sigma_t = \sigma_v/(\kappa k_e p\,\omega_1^2)$
        y $\sigma_n/n = \sqrt2\,\sigma_v/(\pi\kappa k_e\omega_1)$.
  \item \textbf{Retardo del filtro RC del detector:}
        $\varphi_f = \arctan(\omega_e R_f C_f)$ atrasa la conmutación a alta
        velocidad.
  \item \textbf{\emph{Blanking} postconmutación:} ignorar el comparador durante
        $t_b$ (\texttt{BEMF\_BLANKING\_US}, default $0\,\mu\text{s}$, máximo
        $200\,\mu\text{s}$) limita la velocidad máxima a
        $\omega_{m,\text{blank}} = \pi/(6 p t_b)$ ($\approx \SI{12500}{rpm}$ con
        $p=2, t_b=\SI{200}{\micro s}$), muy por encima del rango del agitador.
  \item \textbf{Consenso trifásico:} \texttt{calculate\_consensus\_speed()}
        exige al menos 2 fases consistentes dentro del \SI{15}{\percent} entre sí
        (o degrada a 1 fase válida antes de activar el fallback).
\end{itemize}

\section{Implementación en punto fijo Q16.16 del PI en el STM32}\label{ap:q16}
Para evitar aritmética en coma flotante en el STM32F103 (sin FPU),
\texttt{rpm\_pi\_controller.c} y su réplica exacta en MATLAB
\texttt{firmware\_pi\_reference\_step.m} ejecutan~\eqref{eq:pi_tustin}--\eqref{eq:antiwindup}
en enteros con escala $2^{16}$ (Q16.16):
\begin{itemize}
  \item $K_{p,q} = \mathrm{round}(K_p\cdot 2^{16})$ y $e_k\in\mathbb Z$
        (cuentas de rpm) producen $K_{p,q}e_k$ en Q16;
  \item $(K_i\Delta t_c/2)$ se almacena en Q30 ($K_{i,q30}=\mathrm{round}(K_i\frac{\Delta t_c}{2}2^{30})$)
        para conservar resolución, y el producto por $(e_k+e_{k-1})$ se desplaza
        14 bits a la derecha para volver a Q16 ($30-14=16$);
  \item el \emph{clamping} del integrador y de la salida se aplican en Q16 entre
        $100\cdot 2^{16}$ y $2000\cdot 2^{16}$;
  \item finalmente $u_k = \lfloor u_q/2^{16}\rfloor$ y se escala al período real
        del timer PWM mediante $\mathrm{PWM}=\lfloor(u_k\,\mathrm{ARR}+1000)/2000\rfloor$.
\end{itemize}
Además, las ganancias viajan por el protocolo binario en centésimas de
resolución ($0{,}01$).

% =====================================================================
\begin{thebibliography}{99}
% =====================================================================

\bibitem{pillay1989}
P.~Pillay y R.~Krishnan,
``Modeling, simulation, and analysis of permanent-magnet motor drives,
Part~II: The brushless DC motor drive,''
\emph{IEEE Transactions on Industry Applications}, vol.~25, no.~2,
pp.~274--279, 1989.

\bibitem{hanselman2006}
D.~Hanselman,
\emph{Brushless Permanent Magnet Motor Design}, 2.ª ed.
Magna Physics Publishing, 2006.

\bibitem{griffiths2017}
D.~J.~Griffiths,
\emph{Introduction to Electrodynamics}, 4.ª ed.
Cambridge University Press, 2017. (Cap.~6: torque y energía de un dipolo.)

\bibitem{ravaud2009}
R.~Ravaud, G.~Lemarquand, V.~Lemarquand y C.~Depollier,
``Torque in permanent magnet couplings: Comparison of uniform and radial
magnetization,''
\emph{Journal of Applied Physics}, vol.~105, 053904, 2009.

\bibitem{montague2012}
R.~G.~Montague, C.~M.~Bingham y K.~Atallah,
``Servo control of magnetic gears,''
\emph{IEEE/ASME Transactions on Mechatronics}, vol.~17, no.~2,
pp.~269--278, 2012.

\bibitem{armstrong1994}
B.~Armstrong-Hélouvry, P.~Dupont y C.~Canudas de Wit,
``A survey of models, analysis tools and compensation methods for the control
of machines with friction,''
\emph{Automatica}, vol.~30, no.~7, pp.~1083--1138, 1994.

\bibitem{canudas1995}
C.~Canudas de Wit, H.~Olsson, K.~J.~Åström y P.~Lischinsky,
``A new model for control of systems with friction,''
\emph{IEEE Transactions on Automatic Control}, vol.~40, no.~3,
pp.~419--425, 1995.

\bibitem{simscape-friction}
The MathWorks,
``Rotational Friction,'' documentación de Simscape.\\
\url{https://www.mathworks.com/help/simscape/ref/rotationalfriction.html}

\bibitem{buckingham1914}
E.~Buckingham,
``On physically similar systems; illustrations of the use of dimensional
equations,''
\emph{Physical Review}, vol.~4, no.~4, pp.~345--376, 1914.

\bibitem{rushton1950}
J.~H.~Rushton, E.~W.~Costich y H.~J.~Everett,
``Power characteristics of mixing impellers,''
\emph{Chemical Engineering Progress}, vol.~46, pp.~395--404 y 467--476, 1950.

\bibitem{paul2004}
E.~L.~Paul, V.~A.~Atiemo-Obeng y S.~M.~Kresta (eds.),
\emph{Handbook of Industrial Mixing: Science and Practice}.
Wiley-Interscience, 2004.

\bibitem{mccabe2005}
W.~L.~McCabe, J.~C.~Smith y P.~Harriott,
\emph{Unit Operations of Chemical Engineering}, 7.ª ed.
McGraw-Hill, 2005. (Cap.~9: agitación y mezcla.)

\bibitem{ohmae1982}
T.~Ohmae, T.~Matsuda, K.~Kamiyama y M.~Tachikawa,
``A microprocessor-controlled high-accuracy wide-range speed regulator for
motor drives,''
\emph{IEEE Transactions on Industrial Electronics}, vol.~IE-29, no.~3,
pp.~207--211, 1982.

\end{thebibliography}

\end{document}
